%%%%%%%%%%%%%%%%%%%%%%%%%%%%%%%%%%%%%%%%%%%%%%%%%%%%%%%%%%%%%%%%%%%%%%
% Lớp 10 | Thống kê | Các số đặc trưng đo mức độ phân tán của mẫu số liệu
% Dạng 0D6?4-2: Khoảng biến thiên, khoảng tứ phân vị 
% Số câu: 7  (sinh tự động bởi tach_cau.py; từ nay thầy sửa trực tiếp file này)
%%%%%%%%%%%%%%%%%%%%%%%%%%%%%%%%%%%%%%%%%%%%%%%%%%%%%%%%%%%%%%%%%%%%%%

%%%%%%%%%%%%%%%% TRẮC NGHIỆM 4 PHƯƠNG ÁN %%%%%%%%%%%%%%%%

% nguồn: MaHoaNH_TracNghiem_Lop10
\begin{ex}%[0D6N4-2]%[ma:e65b7a62]
	Khoảng tứ phân vị của dãy  số $2$; $3$; $4$; $5$; $6$ là
	\choice 
		{\True  $\Delta_Q=3$}
		 {$\Delta_Q=\sqrt{2}$}
		{$\Delta_Q=2$} 
		 {$\Delta_Q=-2$}
	 \loigiai{
		Ta  có \[\Delta_Q=Q_3-Q_1=\dfrac{11}{2}-\dfrac{5}{2}=3.\]
	 }
\end{ex}

% nguồn: MaHoaNH_TracNghiem_Lop10
\begin{ex}%[0D6N4-2]%[ma:a76df3f2]
	Cho dãy số liệu thống kê $1$; $2$; $3$; $4$; $5$; $6$; $7$. Khoảng biến  thiên là
	\choice 
		 {$1$}
		 {$2$}
		{$3$} 
		{\True  $6$}
	\loigiai{ 
		Ta có $7-1=6$. 
	} 
\end{ex}

% nguồn: MaHoaNH_TracNghiem_Lop10
\begin{ex}%[0D6N4-2]%[ma:ec3722ab]
	Cho mẫu số liệu thống kê $\{ 1;2;3;4;5;6;7;8;9 \}$. Tìm khoảng tứ phân  vị của mẫu số liệu trên.
	 \choice
		 {$2$}
		 {\True $5$}
		 {$3$}
		{$4$} 
	 \loigiai{
		Ta có $Q_1=2{,}5$; $Q_3=7{,}5$ nên suy ra \[\Delta_Q=Q_3-Q_1=5.\] 
	 }
\end{ex}

% nguồn: MaHoaNH_TracNghiem_Lop10
\begin{ex}%[0D6N4-2]%[ma:69d41c40]
	Cho  mẫu số liệu
	$$\begin{array}{lllllllll}
	3 & 6 & 7 & 8 &  9 & 10 & 12 & 16 & 19
\end{array}$$
Khoảng  tứ phân vị của mẫu số liệu trên là
 \choice
	{$9$} 
	{$6{,}5$} 
	 {\True $7{,}5$}
	 {$14$}
\loigiai{ 
	 Ta có $Q_2=9$; $Q_1=\dfrac{6+7}{2}=6{,}5$; $Q_3=\dfrac{12+16}{2}=14$.\\
	 Khi đó khoảng tứ phân vị $\Delta_Q=Q_3-Q_1=14-6{,}5=7{,}5$.
} 
\end{ex}

% nguồn: MaHoaNH_TracNghiem_Lop10
\begin{ex}%[0D6N4-2]%[ma:7f5573c1]
	Mẫu số liệu sau cho biết cân nặng  (đơn vị kg) của $7$ học sinh lớp $10A$
	$$ 
	\begin{array}{lllllll}
		45 & 56 & 50 & 44 &  42 & 40 & 60
	\end{array}
	$$ 
	Khoảng biến thiên của mẫu số liệu  trên là
	 \choice
		{$22$} 
		 {\True $20$}
		{$42$} 
		{$38$} 
	\loigiai{ 
		Khoảng  biến thiên là $R=x_\text{max}-x_\text{min}=60-40=20$.
	} 
\end{ex}

% nguồn: MaHoaNH_TracNghiem_Lop10
\begin{ex}%[0D6N4-2]%[ma:084bd04e]
	Cho mẫu số liệu sau 
	$$\begin{array}{lllllllll}
	158 & 146 & 152 & 160 & 170 & 154 & 175 & 155 & x 
\end{array}$$
Khi đó $x$ nhận giá  trị nào sau đây đề mẫu số liệu này có khoảng biến thiên là $30$?
 \choice
	 {$160$}
	 {$130$}
	 {\True $145$}
	{$180$} 
 \loigiai{
	\begin{itemize}
		\item  Với $x=160$ khi đó khoảng biên thiên là $175-146=29$.
		 \item Với $x=130$ khi đó khoảng biên thiên là $175-130=45$.
		\item  Với $x=145$ khi đó khoảng biên thiên là $175-145=30$.
		\item Với $x=180$ khi đó khoảng biên  thiên là $180-146=34$.
	\end{itemize}
	Vậy $x=145$. 
} 
\end{ex}


%%%%%%%%%%%%%%%% ĐÚNG/SAI %%%%%%%%%%%%%%%%

% nguồn: MaHoaNH_DungSai_Lop10
\begin{ex}%[0D6H4-2]%[ma:67de4ccd]
	Cho mẫu số liệu $15$; $20$;  $1$; $2$; $4$; $3$; $7$; $5$.
	\\ Các mệnh đề sau  đúng hay sai?
	 \choiceTF
		{\True Cỡ mẫu  là $8$}
		{Khoảng biến thiên của mẫu số liệu  là $10$}
		{\True Mẫu  số liệu có giá trị trung vị là $4{,}5$}
		{Khoảng tứ phân vị của mẫu là  $8$}
	 \loigiai{
		Sắp xếp mẫu số liệu theo thứ tự không giảm  ta được
		 \[1;2;3;4;5;7;15;20.\]
		\begin{itemchoice}
			\itemch Đúng. Từ mẫu số liệu ta có cỡ mẫu là $8$. 
			 \itemch Sai. Khoảng biến thiên là \[20-1=19.\]
			\itemch  Đúng. Cỡ mẫu $n=8$ là số chẵn nên trung vị là \[M_e=\dfrac{1}{2}(4+5)=4{,}5.\]
			\itemch Sai. Tứ phân vị thứ nhất và  thứ ba của mẫu lần lượt là
			\[Q_1=\dfrac{1}{2}(2+3)=2{,}5;  Q_3=\dfrac{1}{2}(7+15)=11.\]
			Suy ra khoảng tứ  phân vị của mẫu là
			\[\Delta_Q=Q_3-Q_1=8{,}5.\] 
		\end{itemchoice}
	} 
\end{ex}
